\documentclass[10pt,a4paper]{amsart}
\usepackage{amsmath,amssymb,amsthm,graphicx}
\usepackage{verbatim}
\newenvironment{reference}{\comment}{\endcomment}
%\newenvironment{reference}{}{}
\newenvironment{slogan}{\comment}{\endcomment}
\newenvironment{history}{\comment}{\endcomment}

% For commutative diagrams we use Xy-pic
\usepackage[all]{xy}

% We use 2cell for 2-commutative diagrams.
\xyoption{2cell}
% Theorem environments.
%
\theoremstyle{plain}
\newtheorem{theorem}[subsection]{Theorem}
\newtheorem{proposition}[subsection]{Proposition}
\newtheorem{lemma}[subsection]{Lemma}

\theoremstyle{definition}
\newtheorem{definition}[subsection]{Definition}
\newtheorem{example}[subsection]{Example}
\newtheorem{exercise}[subsection]{Exercise}
\newtheorem{situation}[subsection]{Situation}

\theoremstyle{remark}
\newtheorem{remark}[subsection]{Remark}
\newtheorem{remarks}[subsection]{Remarks}

\numberwithin{equation}{subsection}

% Macros
%
\def\lim{\mathop{\mathrm{lim}}\nolimits}
\def\colim{\mathop{\mathrm{colim}}\nolimits}
\def\Spec{\mathop{\mathrm{Spec}}}
\def\Hom{\mathop{\mathrm{Hom}}\nolimits}
\def\Ext{\mathop{\mathrm{Ext}}\nolimits}
\def\SheafHom{\mathop{\mathcal{H}\!\mathit{om}}\nolimits}
\def\SheafExt{\mathop{\mathcal{E}\!\mathit{xt}}\nolimits}
\def\Sch{\mathit{Sch}}
\def\Mor{\mathop{\mathrm{Mor}}\nolimits}
\def\Ob{\mathop{\mathrm{Ob}}\nolimits}
\def\Sh{\mathop{\mathit{Sh}}\nolimits}
\def\NL{\mathop{N\!L}\nolimits}
\def\CH{\mathop{\mathrm{CH}}\nolimits}
\def\proetale{{pro\text{-}\acute{e}tale}}
\def\etale{{\acute{e}tale}}
\def\QCoh{\mathit{QCoh}}
\def\Ker{\mathop{\mathrm{Ker}}}
\def\Im{\mathop{\mathrm{Im}}}
\def\Coker{\mathop{\mathrm{Coker}}}
\def\Coim{\mathop{\mathrm{Coim}}}

% Boxtimes
%
\DeclareMathSymbol{\boxtimes}{\mathbin}{AMSa}{"02}

%
% Macros for moduli stacks/spaces
%
\def\QCohstack{\mathcal{QC}\!\mathit{oh}}
\def\Cohstack{\mathcal{C}\!\mathit{oh}}
\def\Spacesstack{\mathcal{S}\!\mathit{paces}}
\def\Quotfunctor{\mathrm{Quot}}
\def\Hilbfunctor{\mathrm{Hilb}}
\def\Curvesstack{\mathcal{C}\!\mathit{urves}}
\def\Polarizedstack{\mathcal{P}\!\mathit{olarized}}
\def\Complexesstack{\mathcal{C}\!\mathit{omplexes}}
% \Pic is the operator that assigns to X its picard group, usage \Pic(X)
% \Picardstack_{X/B} denotes the Picard stack of X over B
% \Picardfunctor_{X/B} denotes the Picard functor of X over B
\def\Pic{\mathop{\mathrm{Pic}}\nolimits}
\def\Picardstack{\mathcal{P}\!\mathit{ic}}
\def\Picardfunctor{\mathrm{Pic}}
\def\Deformationcategory{\mathcal{D}\!\mathit{ef}}
\usepackage{cleveref}
\crefname{originalreference}{original reference}{original references}
\Crefname{originalreference}{Original reference}{Original references}
\makeatletter
\newlabel{situation-shriek}{{0F42}{source}{}{sourcepointer-1}{}}
\newlabel{situation-shriek@cref}{{[originalreference][0][]0F42}{[1][0][]source}{}{}{}}
\makeatother
\crefname{originallemma}{original lemma}{original lemmas}
\Crefname{originallemma}{Original lemma}{Original lemmas}
\makeatletter
\newlabel{lemma-shriek-well-defined}{{0AA0}{source}{}{sourcepointer-2}{}}
\newlabel{lemma-shriek-well-defined@cref}{{[originallemma][0][]0AA0}{[1][0][]source}{}{}{}}
\makeatother
\crefname{originallemma}{original lemma}{original lemmas}
\Crefname{originallemma}{Original lemma}{Original lemmas}
\makeatletter
\newlabel{lemma-upper-shriek-composition}{{0ATX}{source}{}{sourcepointer-3}{}}
\newlabel{lemma-upper-shriek-composition@cref}{{[originallemma][0][]0ATX}{[1][0][]source}{}{}{}}
\makeatother
\crefname{originaldefinition}{original definition}{original definitions}
\Crefname{originaldefinition}{Original definition}{Original definitions}
\makeatletter
\newlabel{categories-definition-functor-into-2-category}{{003N}{source}{}{sourcepointer-4}{}}
\newlabel{categories-definition-functor-into-2-category@cref}{{[originaldefinition][0][]003N}{[1][0][]source}{}{}{}}
\makeatother
\crefname{originaltheorem}{original theorem}{original theorems}
\Crefname{originaltheorem}{Original theorem}{Original theorems}
\makeatletter
\newlabel{flat-theorem-nagata}{{0F41}{source}{}{sourcepointer-5}{}}
\newlabel{flat-theorem-nagata@cref}{{[originaltheorem][0][]0F41}{[1][0][]source}{}{}{}}
\makeatother
\crefname{originallemma}{original lemma}{original lemmas}
\Crefname{originallemma}{Original lemma}{Original lemmas}
\makeatletter
\newlabel{flat-lemma-compactifyable}{{0A9Z}{source}{}{sourcepointer-6}{}}
\newlabel{flat-lemma-compactifyable@cref}{{[originallemma][0][]0A9Z}{[1][0][]source}{}{}{}}
\makeatother
\crefname{originalequation}{original equation}{original equations}
\Crefname{originalequation}{Original equation}{Original equations}
\makeatletter
\newlabel{equation-sheafy}{{0A9L}{source}{}{sourcepointer-7}{}}
\newlabel{equation-sheafy@cref}{{[originalequation][0][]0A9L}{[1][0][]source}{}{}{}}
\makeatother
\crefname{originallemma}{original lemma}{original lemmas}
\Crefname{originallemma}{Original lemma}{Original lemmas}
\makeatletter
\newlabel{lemma-compose-base-change-maps}{{0ATQ}{source}{}{sourcepointer-8}{}}
\newlabel{lemma-compose-base-change-maps@cref}{{[originallemma][0][]0ATQ}{[1][0][]source}{}{}{}}
\makeatother
\crefname{originallemma}{original lemma}{original lemmas}
\Crefname{originallemma}{Original lemma}{Original lemmas}
\makeatletter
\newlabel{lemma-compose-base-change-maps-horizontal}{{0ATR}{source}{}{sourcepointer-9}{}}
\newlabel{lemma-compose-base-change-maps-horizontal@cref}{{[originallemma][0][]0ATR}{[1][0][]source}{}{}{}}
\makeatother
\crefname{originalreference}{original reference}{original references}
\Crefname{originalreference}{Original reference}{Original references}
\makeatletter
\newlabel{remark-going-around}{{0ATS}{source}{}{sourcepointer-10}{}}
\newlabel{remark-going-around@cref}{{[originalreference][0][]0ATS}{[1][0][]source}{}{}{}}
\makeatother
\title{The upper shriek constructions form a pseudo functor on schemes of finite type over the base}
\author{The Stacks Project Authors}
\date{stacks-project:duality.tex:local-snapshot-20261012; GNU FDL 1.2 or later}
\begin{document}
\maketitle
\noindent\textit{Editorial note.} All delivered mathematical text is the author's own, copied from the pinned Stacks Project chapter file duality.tex; no statement or proof step was rewritten, reordered or removed.\par\medskip
\noindent\textit{Editorial note.} Preamble provenance: the pinned chapter file opens with an input of the shared preamble and carries no preamble text of its own, so the project's single Stacks-wide preamble staged at sources/stacks-duality/preamble.tex is excerpted verbatim under the recipe fixed by LANE\_BRIEF 22.1, as three spans, lines 8-13 (verbatim package and the reference/slogan/history environments), lines 14-18 (Xy-pic setup) and lines 154-217 (the native theorem declarations and mathematical macros). Nothing was added, removed or redefined; the wrapper class and the amsmath/amssymb/amsthm/graphicx base packages are supplied by the assembler.\par\medskip
\noindent\textit{Editorial note.} The counted statement and proof contain no cross reference, no citation and no included file, so the delivered body is self contained.\par\medskip
\noindent\textit{Editorial note.} Canonical location of the counted proposition: Stacks tag 0ATY (sources/stacks-duality/pinned-original-source.tex line 4370).\par\medskip
\noindent\textit{Allowed standard prerequisites (editorial).}\begin{itemize}
\item The arrow (g o f)\textasciicircum{}! -> f\textasciicircum{}! o g\textasciicircum{}! and its well definedness. Same chapter, Lemmas labelled lemma-shriek-well-defined and lemma-upper-shriek-composition, cited at proof lines 4378-4380. Supplies the transformations whose coherence is asserted.
\item The base change map for a cartesian square and the two composition rules. Same chapter, equation labelled equation-sheafy, Lemma labelled lemma-compose-base-change-maps and Lemma labelled lemma-compose-base-change-maps-horizontal, cited at proof lines 4410-4422. Identify the arrows of the square with the gamma maps and compose them.
\item The commuting of distant squares and the definition of a pseudo functor. Same chapter, Remark labelled remark-going-around, cited at proof lines 4420-4424, and Categories, Definition labelled categories-definition-functor-into-2-category, invoked in the statement. Give the last commutation gamma\_A o gamma\_C = gamma\_C o gamma\_A and the meaning of the assertion.
\end{itemize}\medskip
\ifdefined\hypertarget\hypertarget{sourcepointer-1}{}\fi
\noindent\textit{Source reference.} Original reference 0F42: sources/stacks-duality/pinned-original-source.tex, source lines 3953--3953. Same-chapter pointer into the author environment situation at source line 3953; invoked at proof line None of the counted proof. It is retained only as a navigation pointer; every load-bearing step is inside the delivered span.\par\medskip
\ifdefined\hypertarget\hypertarget{sourcepointer-2}{}\fi
\noindent\textit{Source reference.} Original lemma 0AA0: sources/stacks-duality/pinned-original-source.tex, source lines 3986--3986. Same-chapter pointer into the author environment lemma at source line 3986; invoked at proof line None of the counted proof. It is retained only as a navigation pointer; every load-bearing step is inside the delivered span.\par\medskip
\ifdefined\hypertarget\hypertarget{sourcepointer-3}{}\fi
\noindent\textit{Source reference.} Original lemma 0ATX: sources/stacks-duality/pinned-original-source.tex, source lines 4161--4161. Same-chapter pointer into the author environment lemma at source line 4161; invoked at proof line 4411 of the counted proof. It is retained only as a navigation pointer; every load-bearing step is inside the delivered span.\par\medskip
\ifdefined\hypertarget\hypertarget{sourcepointer-4}{}\fi
\noindent\textit{Source reference.} Original definition 003N: sources/stacks-chow/tags.txt, source lines 141--141. Cross-chapter standard prerequisite located in the project tag list staged in the local snapshot; invoked at proof line None of the counted proof. It is retained only as a navigation pointer; every load-bearing step is inside the delivered span.\par\medskip
\ifdefined\hypertarget\hypertarget{sourcepointer-5}{}\fi
\noindent\textit{Source reference.} Original theorem 0F41: sources/stacks-chow/tags.txt, source lines 18886--18886. Cross-chapter standard prerequisite located in the project tag list staged in the local snapshot; invoked at proof line 4394 of the counted proof. It is retained only as a navigation pointer; every load-bearing step is inside the delivered span.\par\medskip
\ifdefined\hypertarget\hypertarget{sourcepointer-6}{}\fi
\noindent\textit{Source reference.} Original lemma 0A9Z: sources/stacks-chow/tags.txt, source lines 12967--12967. Cross-chapter standard prerequisite located in the project tag list staged in the local snapshot; invoked at proof line 4395 of the counted proof. It is retained only as a navigation pointer; every load-bearing step is inside the delivered span.\par\medskip
\ifdefined\hypertarget\hypertarget{sourcepointer-7}{}\fi
\noindent\textit{Source reference.} Original equation 0A9L: sources/stacks-duality/pinned-original-source.tex, source lines 988--988. Same-chapter pointer into the author environment equation at source line 988; invoked at proof line 4413 of the counted proof. It is retained only as a navigation pointer; every load-bearing step is inside the delivered span.\par\medskip
\ifdefined\hypertarget\hypertarget{sourcepointer-8}{}\fi
\noindent\textit{Source reference.} Original lemma 0ATQ: sources/stacks-duality/pinned-original-source.tex, source lines 1238--1238. Same-chapter pointer into the author environment lemma at source line 1238; invoked at proof line 4415 of the counted proof. It is retained only as a navigation pointer; every load-bearing step is inside the delivered span.\par\medskip
\ifdefined\hypertarget\hypertarget{sourcepointer-9}{}\fi
\noindent\textit{Source reference.} Original lemma 0ATR: sources/stacks-duality/pinned-original-source.tex, source lines 1375--1375. Same-chapter pointer into the author environment lemma at source line 1375; invoked at proof line 4416 of the counted proof. It is retained only as a navigation pointer; every load-bearing step is inside the delivered span.\par\medskip
\ifdefined\hypertarget\hypertarget{sourcepointer-10}{}\fi
\noindent\textit{Source reference.} Original reference 0ATS: sources/stacks-duality/pinned-original-source.tex, source lines 1471--1471. Same-chapter pointer into the author environment remark at source line 1471; invoked at proof line 4423 of the counted proof. It is retained only as a navigation pointer; every load-bearing step is inside the delivered span.\par\medskip
\begin{lemma}
\label{lemma-pseudo-functor}
In Situation \ref{situation-shriek} the constructions of
Lemmas \ref{lemma-shriek-well-defined} and \ref{lemma-upper-shriek-composition}
define a pseudo functor from the category $\textit{FTS}_S$
into the $2$-category of categories (see Categories, Definition
\ref{categories-definition-functor-into-2-category}).
\end{lemma}

\begin{proof}
To show this we have to prove given morphisms
$f : X \to Y$, $g : Y \to Z$, $h : Z \to T$
that
$$
\xymatrix{
(h \circ g \circ f)^! \ar[r]_{\gamma_{A + B}} \ar[d]_{\gamma_{B + C}} &
f^! \circ (h \circ g)^! \ar[d]^{\gamma_C} \\
(g \circ f)^! \circ h^! \ar[r]^{\gamma_A} & f^! \circ g^! \circ h^!
}
$$
is commutative (for the meaning of the $\gamma$'s, see below).
To do this we choose a compactification $\overline{Z}$
of $Z$ over $T$, then a compactification $\overline{Y}$ of $Y$ over
$\overline{Z}$, and then a compactification $\overline{X}$ of
$X$ over $\overline{Y}$. This uses
More on Flatness, Theorem \ref{flat-theorem-nagata} and
Lemma \ref{flat-lemma-compactifyable}.
Let $W \subset \overline{Y}$ be the inverse image of $Z$ under
$\overline{Y} \to \overline{Z}$ and let $U \subset V \subset \overline{X}$
be the inverse images of $Y \subset W$ under $\overline{X} \to \overline{Y}$.
This produces the following diagram
$$
\xymatrix{
X \ar[d]_f \ar[r] & U \ar[r] \ar[d] \ar@{}[dr]|A &
V \ar[d] \ar[r] \ar@{}[rd]|B & \overline{X} \ar[d] \\
Y \ar[d]_g \ar[r] & Y \ar[r] \ar[d] & W \ar[r] \ar[d] \ar@{}[rd]|C &
\overline{Y} \ar[d] \\
Z \ar[d]_h \ar[r] & Z \ar[d] \ar[r] & Z \ar[d] \ar[r] & \overline{Z} \ar[d] \\
T \ar[r] & T \ar[r] & T \ar[r] & T
}
$$
Without introducing tons of notation but arguing exactly
as in the proof of Lemma \ref{lemma-upper-shriek-composition}
we see that the maps in the first displayed diagram use the
maps (\ref{equation-sheafy}) for the rectangles
$A + B$, $B + C$, $A$, and $C$ as indicated. Since by
Lemmas \ref{lemma-compose-base-change-maps} and
\ref{lemma-compose-base-change-maps-horizontal}
we have $\gamma_{A + B} = \gamma_A \circ \gamma_B$ and
$\gamma_{B + C} = \gamma_C \circ \gamma_B$  we conclude
that the desired equality holds provided
$\gamma_A \circ \gamma_C = \gamma_C \circ \gamma_A$.
This is true because the two squares $A$ and $C$ only
intersect in one point (similar to the last argument in
Remark \ref{remark-going-around}).
\end{proof}

\end{document}
